\subsection{How we got to the final model}
\label{sec:journey}
The final model was not designed in one step. \Cref{tab:journey} tells the project in the order it happened: each row is something we tried, what it gave, and what that taught us, which is what decided the next row. The complete list of runs with their settings is in \cref{app:log,app:who}.

\begin{table}[p]
\centering\small
\caption{\textbf{From first try to final model.} Accuracies are validation accuracies unless marked $^\dagger$ (private leaderboard). \ver\ verified by the codebase, \rep\ from Erdem's log, the experiment records in the repository or Kaggle; $^\S$ on the DINOv2 branch's development split.}
\label{tab:journey}
\renewcommand{\arraystretch}{1.18}
\newcommand{\stp}[1]{\tikz[baseline=(n.base)]\node[circle, fill=accent, text=white, inner sep=1.2pt, minimum size=4.2mm, font=\sffamily\scriptsize\bfseries](n){#1};}
\begin{tabularx}{\textwidth}{@{}c>{\raggedright\arraybackslash}p{44mm}>{\raggedleft\arraybackslash}p{17mm}>{\raggedright\arraybackslash}X@{}}
\toprule
 & What we tried & Result & What we learned, and what we changed\\
\midrule
\stp{1} & First CNN: 3 conv blocks, no normalisation, Adam at a constant $10^{-3}$, 20 epochs (30 Sept.) & 76.71\%\ver & Training accuracy only 75\%: the model underfits. \textbf{Next:} BatchNorm, a fourth stage and a OneCycle schedule.\\
\stp{2} & EmojiCNNv2 (0.39\,M parameters), same data & 90.96\%\ver & +14 points from optimisation and normalisation alone, no extra data. A good CNN from scratch is a strong baseline.\\
\stp{3} & File metadata only (size, PNG mode, hidden colour), random forest & 87.37\%\ver & The file format leaks the label. \textbf{Next:} composite on white, square canvas, exact-pixel augmentation; keep the final model pixel-only. Same CNN with compositing: 92.11\% $\to$ 93.62\%\rep.\\
\stp{4} & Hand-crafted features: colour histogram + HOG with an SVM (Erdem); extra-trees on pixels, colour and transparency (Lina) & 81.28\%, 82.14\%\rep & Below every trained CNN: the cues are too fine for generic descriptors. \textbf{Next:} learn the features.\\
\stp{5} & Pretrained ResNet18 as a frozen feature extractor: linear head (Erdem); with hand-crafted style features and an RBF SVM (Lina) & 77.48\%, 96.76\%\rep & ImageNet features describe objects, not drawing style; frozen, they need hand-crafted style descriptors to compete. \textbf{Next:} fine-tune the whole network.\\
\stp{6} & Full fine-tuning at 128\,px: ResNet18, EfficientNet-B0, ConvNeXt-Tiny; ResNet18 at 160\,px with file-format features (Charles) & 93.72\%, 93.12\%, 95.90\%\rep; 0.96396$^\dagger$ & Fine-tuning recovers 16 points; ConvNeXt is best. Charles's entry also reads the file format, so it is not a pixels-only result. \textbf{Next:} build on ConvNeXt-Tiny.\\
\stp{7} & Stronger augmentation: rotation, scaling, label smoothing & 94.38\%\rep & $-1.5$ points: resampling blurs the outlines that identify a vendor. \textbf{Next:} flips and integer shifts only.\\
\stp{8} & Input resolution 128 $\to$ 224\,px & 97.57\%\rep & +1.7 points, the largest single gain: resolution matters more than architecture.\\
\stp{9} & Ensembles (128 + 224\,px), retraining on all data, nearest-neighbour upscaling (E15, Erdem) & 97.93\%\rep; 0.98238$^\dagger$ & Retraining on all data changed the private score by +0.14 and $-0.08$ points: no reliable gain. The best entry mixes resolutions and resize filters; in a controlled test at 144\,px, nearest-neighbour instead of bicubic resizing gave +0.9 points (E12a, E12b\rep); at 216\,px nearest-neighbour did no better than 224\,px bicubic (E13, E8\rep).\\
\stp{10} & Self-supervised ViT: DINOv2 ViT-S/14, last 2 of 12 blocks tuned (Lina); averaged with a partially tuned ConvNeXt & 95.90\%, 97.12\%$^\S$; 0.95607$^\dagger$, 0.96477$^\dagger$ & The transformer is weaker here, and the average fell 0.6 points from its development estimate to the leaderboard (reused development split). \textbf{Next:} try learning the combination.\\
\stp{11} & Learned fusion of DINOv2, ConvNeXt and 1{,}460 style features (transformer and MLP heads); decision models Julia-1 and Jev & best 97.01\%$^\S$ vs 97.12\% & No combiner beat the plain average; Jev fixed 8 errors and made 37. \textbf{Next:} back to one strong architecture, fine-tuned end to end.\\
\stp{12} & ConvNeXt-Tiny, 5 grouped folds, flip TTA, lowest-validation-loss epoch (run A) & 97.29\%\ver & Validation loss and accuracy disagree. \textbf{Next:} keep the last epoch of the schedule.\\
\rowcolor{panel}\stp{13} & Same, last epoch (run B): \textbf{final model} & 97.69\%\ver; 0.98218$^\dagger$ & +0.40 points over run A ($p=2.4\times10^{-5}$), 13 minutes of training.\\
\bottomrule
\end{tabularx}
\end{table}

Three turns shaped the project. The first was finding the file-format shortcut (step 3): it changed the question from ``how high can we score'' to ``how high can we score from the drawing'', and every later comparison was made under that rule. The second was the failure of frozen features (step 5), which showed that vendor style is not what ImageNet pretraining encodes, and that the backbone has to be adapted. The third was the gap between validation loss and accuracy at the end (step 12), which changed how we choose a checkpoint. Several ideas did not work, and they are as informative as the ones that did: resampling augmentation, retraining on all data, adding a weaker transformer to a strong CNN, and learned fusion or decision models on top of two networks and hand-crafted style features all lost or failed to gain accuracy (\cref{sec:nowork,sec:fusion}); only an expert that reads the file bytes, combined by simple rules, added accuracy.

\subsection{Experiment ladder}
\label{sec:ladder}
